\section{The formation}
\label{sec:formation}

\subsection{The series circuit}

Heat leaves the material through four resistances in series, per unit length
of borehole:
\begin{equation}
\Rtot(t) = \underbrace{R_p + R_{\mathrm{cas}} + R_{\mathrm{cem}}}_{\Rnear}
  + \Rg(t),
\end{equation}
with
\begin{align}
R_p &= \frac{\ln(R_c/r_e)}{2\pi\,\nft\,k_p},
&
R_{\mathrm{cas}} &= \frac{\ln(r_{co}/R_c)}{2\pi k_{\mathrm{steel}}},
\\[4pt]
R_{\mathrm{cem}} &= \frac{\ln(r_b/r_{co})}{2\pi k_{\mathrm{cem}}},
&
\Rg(t) &= \frac{1}{2\pi k_r}
  \ln\!\left(\frac{2.5\sqrt{\alpha t}}{r_b}\right).
\label{eq:Rg}
\end{align}

\begin{table}[h]
\centering
\begin{tabular}{lrr}
\toprule
term & value, \si{\metre\kelvin\per\watt} & share \\
\midrule
material to casing, $R_p$      & 0.14543 & \SI{26.5}{\percent} \\
casing steel, $R_{\mathrm{cas}}$ & 0.00034 & \SI{0.1}{\percent} \\
cement, $R_{\mathrm{cem}}$     & 0.01838 & \SI{3.4}{\percent} \\
\textbf{ground}, $\Rg$         & \textbf{0.38413} & \textbf{\SI{70.1}{\percent}} \\
\midrule
total                          & 0.54829 & \\
\bottomrule
\end{tabular}
\caption{The resistance budget at the design point, $t_{\mathrm{op}} =
\SI{10}{yr}$. The casing is thermally invisible and the cement is a rounding
error; the ground resistance is the model.}
\end{table}

Two remarks. First, $R_p$ is the crude term: it treats the $\nft$ tubes as
independent parallel paths to the casing and ignores that those paths crowd
together near the wall. It is a quarter of a total the ground dominates, and
refining it would move the result less than the uncertainty in $k_r$ does ---
across the plausible range $k_r \in [1.5, 3.5]$ the loss varies by a factor of
2.1. Second, $\Rg$ is the only term that depends on time. It grows as the
thermal front spreads, which is why $t_{\mathrm{op}}$ is a parameter of the
case and why a quoted loss figure is meaningless without the age beside it.

\subsection{The undisturbed temperature profile}

\begin{equation}
\Trock(z) = T_s + \nabla T \cdot \text{depth}(z).
\end{equation}
At the default \SI{45}{\kelvin\per\kilo\metre} from \SI{18}{\degC}, the rock
runs from \SI{18}{\degC} at the surface to \SI{86}{\degC} at
\SI{1524}{\metre}. At the central Wilmington value of
\SI{56.5}{\kelvin\per\kilo\metre}, compiled in 1949 from 84 temperature
measurements and the highest known in the Los Angeles basin, it reaches
\SI{104}{\degC} at \SI{5000}{ft} and \SI{156}{\degC} at \SI{8000}{ft} ---
\emph{hotter than the coldest cascade layers}. Those segments gain heat from
the formation. The implementation takes no absolute values anywhere, so a
negative loss is reported as a result rather than absorbed.

\subsection{The gradient and the source temperature are independent}

The wellhead temperature of the geothermal stream is a \emph{lower bound} on
the formation temperature at producing depth: the water cools as it ascends
the production string, and the producing wells may be deeper than the
retrofitted ones. Deriving the gradient from the source temperature writes
that bound as an equality and discards precisely the physics the deficit
represents. Both are therefore specified, and only the inequality
\begin{equation}
T_{\mathrm{source}} \;\le\; T_s + \nabla T\, L_{\mathrm{prod}}
\end{equation}
is enforced --- as a check that can fail, with a warning when the margin falls
below \SI{5}{\kelvin}, which would leave no allowance for the ascent.

\subsection{Coupling to the enthalpy march}

Each segment now has two linear sinks acting on it: the tube at $T_0$ with
conductance $K$, and the formation at $\Trock$ with conductance
$K_\ell = 1/(\Rtot \nft)$, the division by $\nft$ because the per-borehole
conductance is shared among the tubes at that depth. These combine exactly:
for any material temperature $T$,
\begin{equation}
K(T_0 - T) + K_\ell(\Trock - T)
  = (K + K_\ell)\,(T_{\mathrm{eq}} - T),
\qquad
T_{\mathrm{eq}} = \frac{K T_0 + K_\ell \Trock}{K + K_\ell}.
\label{eq:combine}
\end{equation}
The segment advance is therefore called with $(K + K_\ell, T_{\mathrm{eq}})$
and is unchanged in form, remaining implicit in $T$. Adding the loss
explicitly after the step would have been a half-order scheme on a term that
reverses sign with depth.

The split is recovered afterwards without further assumption. If $q_{\Sigma}$
is the net heat the material received over the step, the effective material
temperature follows from $q_{\Sigma} = (K+K_\ell)(T_{\mathrm{eq}} - T_{\rm
eff})$, and then
\begin{equation}
q_{\mathrm{loss}} = K_\ell\,(T_{\mathrm{eff}} - \Trock),
\qquad
\qp_{\mathrm{tube}} = q_{\Sigma} + q_{\mathrm{loss}} .
\label{eq:split}
\end{equation}
Only $\qp_{\mathrm{tube}}$ is charged against the fluid: the formation term is
not the fluid's to pay.

\subsection{What the energy closure becomes}

In THUMS the closure was $|Q_{\mathrm{cum}} - \Delta E|/|Q_{\mathrm{cum}}|$:
everything the tube gave went into storage, because nothing could leave. The
identity is now
\begin{equation}
Q_{\mathrm{tube}} = \Delta E + Q_{\mathrm{loss}},
\end{equation}
and the residual is what remains. This is a bookkeeping identity rather than a
gate --- it cannot fail without an arithmetic error --- but it is retained
because it catches one specific mistake that the split in~\eqref{eq:split}
invites, namely omitting $q_{\mathrm{loss}}$ from the tube duty. It does
\emph{not} catch a sign error, for reasons developed in
Section~\ref{sec:gates}.

